% sec2_6.tex — Section 2.6: Implications of Conditional Homoskedasticity

The test statistics developed in Sections~2.4 and 2.5 are different from the
finite-sample counterparts of Section~1.4 designed for testing the same null hypothesis.
How are they related?  What is the asymptotic distribution of the $t$ and $F$ statistics
of Chapter~1?  This section answers these questions.

It turns out that the robust $t$-ratio is numerically equal to the $t$-ratio of
Section~1.4, for a particular choice of $\widehat{\mathbf{S}}$.  Therefore, the
asymptotic distribution of the $t$-ratio of Section~1.4 is the same as that of the
robust $t$-ratio, if that particular choice is consistent for $\mathbf{S}$.  The same
relationship holds between the $F$-ratio of Section~1.4 and the Wald statistic $W$ of
this chapter.  Under the conditional homoskedasticity assumption stated below, that
particular choice is indeed consistent.

%% ─── Conditional versus Unconditional Homoskedasticity ───────────────────────
\subsection*{Conditional versus Unconditional Homoskedasticity}

The conditional homoskedasticity assumption is:

\begin{assumption}[conditional homoskedasticity]
\label{ass:2.7}
\begin{equation}
  \E(\varepsilon_i^2 \mid \mathbf{x}_i) = \sigma^2 > 0.
  \label{eq:2.6.1}
\end{equation}
\end{assumption}

This assumption implies that the unconditional second moment $\E(\varepsilon_i^2)$ equals
$\sigma^2$ by the Law of Total Expectations.  To be clear about the distinction between
unconditional and conditional homoskedasticity, consider the following example.

\begin{example}[unconditionally homoskedastic but conditionally heteroskedastic errors]
\label{ex:2.6}
As already observed, if $\{y_i, \mathbf{x}_i\}$ is stationary, so is
$\{\varepsilon_i\}$, and the error is \emph{un}conditionally homoskedastic in that
$\E(\varepsilon_i^2)$ does not depend on $i$.  To illustrate that the error can
nevertheless be conditionally heteroskedastic, suppose that $\varepsilon_i$ is written as
$\varepsilon_i \equiv \eta_i f(\mathbf{x}_i)$, where $\{\eta_i\}$ is zero-mean
$\E(\eta_i) = 0$ and is independent of $\mathbf{x}_i$.  The conditional second moment of
$\varepsilon_i$ depends on $\mathbf{x}_i$ because
\begin{align*}
  \E(\varepsilon_i^2 \mid \mathbf{x}_i)
    &= \E(\eta_i^2 f(\mathbf{x}_i)^2 \mid \mathbf{x}_i)
    \quad (\text{since } \varepsilon_i \equiv \eta_i f(\mathbf{x}_i)) \\
    &= f(\mathbf{x}_i)^2 \,\E(\eta_i^2 \mid \mathbf{x}_i)
    \quad (\text{by the linearity of conditional expectations}) \\
    &= f(\mathbf{x}_i)^2 \,\E(\eta_i^2)
    \quad (\text{since } \eta_i \text{ is independent of } \mathbf{x}_i
      \text{ by assumption}),
\end{align*}
which varies across $i$ because of the variation in $f(\mathbf{x}_i)$ across $i$.
\end{example}

%% ─── Reduction to Finite-Sample Formulas ─────────────────────────────────────
\subsection*{Reduction to Finite-Sample Formulas}

To examine the large-sample distribution of the $t$ and $F$ statistics of Chapter~1
under the additional Assumption~2.7, we first examine the \emph{algebraic} relationship
to their robust counterparts.  Consider the following choice for the estimate of
$\mathbf{S}$:
\[
  \widehat{\mathbf{S}} = s^2 \mathbf{S}_{\mathbf{xx}},
\]
where $s^2$ is the OLS estimate of $\sigma^2$.  (We will show in a moment that this
estimator is consistent under conditional homoskedasticity.)  Then the expression (2.3.5)
for $\widehat{\avar}(\mathbf{b})$ becomes
\begin{equation}
  \widehat{\avar}(\mathbf{b}) = s^2 \mathbf{S}_{\mathbf{xx}}^{-1}
    = n \cdot s^2 \cdot (\mathbf{X}'\mathbf{X})^{-1}.
  \label{eq:2.6.2}
\end{equation}
Substituting this expression into (2.4.1), we see that the robust standard error becomes
\begin{equation}
  \sqrt{s^2 \text{ times } (k,k) \text{ element of } (\mathbf{X}'\mathbf{X})^{-1}},
  \label{eq:2.6.3}
\end{equation}
which is the usual standard error in finite-sample theory.  So the robust $t$-ratio is
numerically identical to the usual finite-sample $t$-ratio when we set
$\widehat{\mathbf{S}} = s^2 \mathbf{S}_{\mathbf{xx}}$.  Similarly, substituting
\eqref{eq:2.6.2} into the expression for the Wald statistic (2.4.2), we obtain
\begin{align}
  W &= n \cdot (\mathbf{R}\mathbf{b} - \mathbf{r})'
    \bigl\{\mathbf{R}[n \cdot s^2 \cdot (\mathbf{X}'\mathbf{X})^{-1}]\mathbf{R}'
    \bigr\}^{-1} (\mathbf{R}\mathbf{b} - \mathbf{r}) \notag \\
    &= (\mathbf{R}\mathbf{b} - \mathbf{r})'
    \bigl\{\mathbf{R}[s^2 \cdot (\mathbf{X}'\mathbf{X})^{-1}]\mathbf{R}'
    \bigr\}^{-1} (\mathbf{R}\mathbf{b} - \mathbf{r})
    \quad (\text{the two } n\text{'s cancel}) \notag \\
    &= (\mathbf{R}\mathbf{b} - \mathbf{r})'
    \bigl\{\mathbf{R}[(\mathbf{X}'\mathbf{X})^{-1}]\mathbf{R}'
    \bigr\}^{-1} (\mathbf{R}\mathbf{b} - \mathbf{r}) / s^2 \notag \\
    &= r \cdot F \quad (\text{by the definition of (1.4.9) of the } F\text{-ratio})
      \notag \\
    &= (SSR_R - SSR_U)/s^2 \quad (\text{by (1.4.11)}).
\end{align}
Thus, when we set $\widehat{\mathbf{S}} = s^2 \mathbf{S}_{\mathbf{xx}}$, the Wald
statistic $W$ is numerically identical to $\#r \cdot F$ (where $\#r$ is the number of
restrictions in the null hypothesis).

%% ─── Large-Sample Distribution of t and F Statistics ─────────────────────────
\subsection*{Large-Sample Distribution of $t$ and $F$ Statistics}

It then follows from Proposition~2.3 that the $t$-ratio (2.4.1) is asymptotically
$N(0,1)$ and $\#r \cdot F$ asymptotically $\chi^2(\#r)$, if
$s^2 \mathbf{S}_{\mathbf{xx}}$ is consistent for $\mathbf{S}$.  That
$s^2 \mathbf{S}_{\mathbf{xx}}$ is consistent for $\mathbf{S}$ can be seen as follows.
Under conditional homoskedasticity, the matrix of fourth moments $\mathbf{S}$ can be
expressed as a product of second moments:
\begin{align}
  \mathbf{S} &= \E(\mathbf{g}_i \mathbf{g}_i')
    = \E(\mathbf{x}_i \mathbf{x}_i' \varepsilon_i^2)
    \quad (\text{since } \mathbf{g}_i \equiv \mathbf{x}_i \cdot \varepsilon_i) \notag \\
    &= \E[\E(\mathbf{x}_i \mathbf{x}_i' \varepsilon_i^2 \mid \mathbf{x}_i)]
    \quad (\text{by the Law of Total Expectations}) \notag \\
    &= \E[\mathbf{x}_i \mathbf{x}_i' \,\E(\varepsilon_i^2 \mid \mathbf{x}_i)]
    \quad (\text{by the linearity of conditional expectations}) \notag \\
    &= \E(\mathbf{x}_i \mathbf{x}_i' \sigma^2)
    \quad (\text{by Assumption~2.7}) \notag \\
    &= \sigma^2 \,\E(\mathbf{x}_i \mathbf{x}_i')
    = \sigma^2 \bm{\Sigma}_{\mathbf{xx}}.
  \label{eq:2.6.4}
\end{align}
This decomposition has several implications.
\begin{itemize}
  \item ($\bm{\Sigma}_{\mathbf{xx}}$ is nonsingular) \;
    Since by Assumption~2.5 $\mathbf{S}$ is nonsingular, this decomposition of
    $\mathbf{S}$ implies that $\sigma^2 > 0$ and $\bm{\Sigma}_{\mathbf{xx}}$ is
    nonsingular.  Hence, Assumption~2.4 (rank condition) is implied.

  \item (No need for fourth-moment assumption) \;
    By ergodic stationarity $\mathbf{S}_{\mathbf{xx}} \pto \bm{\Sigma}_{\mathbf{xx}}$.
    By Proposition~2.2, $s^2$ is consistent for $\sigma^2$ under Assumptions~2.1--2.4.
    Thus, $s^2 \mathbf{S}_{\mathbf{xx}} \pto \sigma^2 \bm{\Sigma}_{\mathbf{xx}}
    = \mathbf{S}$.  We do not need the fourth-moment assumption (Assumption~2.6) for
    consistency.
\end{itemize}

As another implication of \eqref{eq:2.6.4}, the expression for $\avar(\mathbf{b})$ can
be simplified: inserting \eqref{eq:2.6.4} into (2.3.4) of Proposition~2.1, the expression
for $\avar(\mathbf{b})$ becomes
\begin{equation}
  \avar(\mathbf{b}) = \sigma^2 \bm{\Sigma}_{\mathbf{xx}}^{-1}.
  \label{eq:2.6.5}
\end{equation}
Thus, we have proved

\begin{proposition}[large-sample properties of $\mathbf{b}$, $t$, and $F$ under
conditional homoskedasticity]
\label{prop:2.5}
\emph{Suppose Assumptions~2.1--2.5 and 2.7 are satisfied.  Then}
\begin{enumerate}[label=(\alph*)]
  \item \emph{(Asymptotic distribution of $\mathbf{b}$) \; The OLS estimator
    $\mathbf{b}$ of $\bm{\beta}$ is consistent and asymptotically normal with}
    $\avar(\mathbf{b}) = \sigma^2 \bm{\Sigma}_{\mathbf{xx}}^{-1}$.

  \item \emph{(Consistent estimation of asymptotic variance) \; Under the same set of
    assumptions, $\avar(\mathbf{b})$ is consistently estimated by}
    $\widehat{\avar}(\mathbf{b}) = s^2 \mathbf{S}_{\mathbf{xx}}^{-1}
    = n \cdot s^2 \cdot (\mathbf{X}'\mathbf{X})^{-1}$.

  \item \emph{(Asymptotic distribution of the $t$ and $F$ statistics of the
    finite-sample theory) \; Under $\mathrm{H}_0\colon \beta_k = \bar{\beta}_k$, the
    usual $t$-ratio (1.4.5) is asymptotically distributed as $N(0,1)$.  Under
    $\mathrm{H}_0\colon \mathbf{R}\bm{\beta} = \mathbf{r}$,
    $\#r \cdot F$ is asymptotically $\chi^2(\#r)$, where $F$ is the $F$ statistic from
    (1.4.9) and $\#r$ is the number of restrictions in $\mathrm{H}_0$.}
\end{enumerate}
\end{proposition}

%% ─── Variations of Asymptotic Tests ──────────────────────────────────────────
\subsection*{Variations of Asymptotic Tests under Conditional Homoskedasticity}

According to this result, you should look up the $N(0,1)$ table to find the critical
value to be compared with the $t$-ratio (1.4.5) and the $\chi^2$ table for the statistic
$\#r \cdot F$ derived from (1.4.9).  Some researchers replace the $s^2$ in (1.4.5) and
(1.4.9) by $\frac{1}{n} \sum_i e_i^2$.  That is, the degrees of freedom $n - K$ is
replaced by $n$, or the degrees of freedom adjustment implicit in $s^2$ is removed.
The difference this substitution makes vanishes in large samples, because
$\lim_{n \to \infty} [n/(n - K)] = 1$.  Therefore, regardless of which test to use, the
outcome of the test will be the same if the sample size is sufficiently large.

Another variation is to retain the degrees of freedom $n - K$ but use the $t(n - K)$
table for the $t$ ratio and the $F(\#r, n - K)$ table for $F$, which is \emph{exactly}
the prescription of finite-sample theory.  This, too, is asymptotically valid because, as
$n - K$ tends to infinity (which is what happens when $n \to \infty$ with $K$ fixed),
the $t(n - K)$ distribution converges to $N(0,1)$ (just compare the $t$ table for large
degrees of freedom with the standard normal table) and $F(\#r, n - K)$ to
$\chi^2(\#r)/\#r$.  Put differently, even if the error is not normally distributed and the
regressors are merely predetermined (orthogonal to the error term) and not strictly
exogenous, the distribution of the $t$-ratio (1.4.5) is well approximated by $t(n - K)$,
and the distribution of the $F$ ratio by $F(\#r, n - K)$.

These variations are all asymptotically equivalent in that the differences in values
vanishes in large samples and hence (by Lemma~2.4(a)) their asymptotic distributions are
the same.  However, when the sample size is only moderately large, the approximation to
the finite-sample or exact distribution of test statistics may be better with $t(n - K)$
and $F(\#r, n - K)$, rather than with $N(0,1)$ and $\chi^2(\#r)$.  Because the exact
distribution depends on the DGP, there is no simple guide as to which variation works
better in finite samples.  This issue of which table --- $N(0,1)$
or $t(n - K)$ --- should be used for moderate sample sizes will be taken up in the
Monte Carlo exercise of this chapter.
